\documentclass[a4paper,12pt]{article}

\usepackage[utf8]{inputenc}
\usepackage[T1]{fontenc}
\usepackage[italian]{babel}
\usepackage[margin=2.5cm]{geometry}
\usepackage{graphicx}
\usepackage{grffile}
\usepackage{booktabs}
\usepackage{setspace}
\usepackage{titlesec}
\usepackage{float}
\usepackage{ifthen}
\usepackage{tcolorbox}
\usepackage{enumitem}
\usepackage[colorlinks=true,linkcolor=black,urlcolor=primaryblue,citecolor=primaryblue]{hyperref}
\usepackage{changepage}
\usepackage{xspace}   % <--- AGGIUNTO QUI

\usepackage[table]{xcolor}
\usepackage{tabularx}
\usepackage{array}

% ----- Macro (corrette) -----
\definecolor{primaryblue}{RGB}{0,102,204}
\definecolor{secondaryblue}{RGB}{51,153,255}
\definecolor{lightgray}{RGB}{245,245,245}
\definecolor{darkgray}{RGB}{100,100,100}

\newcommand{\G}{\textsubscript{\scalebox{0.6}{\textbf{G}}}\xspace}

% Comando per la riga delle decisioni (senza \arrayrulecolor interno)
\newcommand{\DecisionRow}[3]{%
    \rowcolor{secondaryblue!10} #1 & #2 & #3 \\
    \hline
}

% Logo con fallback
\newcommand{\VBLogo}[3]{%
    \IfFileExists{#1}{%
        \includegraphics[width=#2,height=#3,keepaspectratio]{#1}%
    }{%
        \fbox{\parbox[c][#3][c]{#2}{\centering Logo non trovato\\(#1)}}%
    }%
}
% ---------------------------

\titleformat{\section}
  {\Large\bfseries\color{primaryblue}}
  {\thesection}{1em}{}

\setlength{\parskip}{4pt}
\setlength{\parindent}{0pt}

\setlist[itemize]{leftmargin=*,itemsep=3pt}
\setlist[enumerate]{leftmargin=*,itemsep=3pt}
\graphicspath{{./}{../assets/images/}{./images/}}

\begin{document}

\begin{center}
  \VBLogo{../../../../assets/Logo.jpg}{6cm}{3cm} \\[0.8cm]

  {\Large\bfseries\color{primaryblue} BugBusters}\\[0.3cm]
  {\small\color{darkgray} Email: \texttt{bugbusters.unipd@gmail.com}} \\[0.1cm]
  {\small\color{darkgray} Gruppo: 4} \\[0.5cm]

  {\large\bfseries Università degli Studi di Padova}\\[0.3cm]
  {\small Laurea in Informatica}\\[0.2cm]
  {\small Corso: Ingegneria del Software}\\[0.2cm]
  {\small Anno Accademico: 2025/2026}\\[0.8cm]

  {\Huge\bfseries\color{primaryblue} Verbale Esterno\G}\\[0.3cm]
  {\Large\color{secondaryblue} 30 marzo 2026}\\[0.8cm]
\end{center}

\begin{center}
\begin{tcolorbox}[colback=lightgray,colframe=primaryblue,width=0.85\textwidth,arc=3mm,boxrule=0.5pt]
\begin{tabularx}{\linewidth}{@{}lX@{}}
\textbf{Redattore\G}    & Alberto Autiero \\
\textbf{Verificatore\G} & Alberto Pignat \\
\textbf{Uso}          & Esterno \\
\textbf{Destinatari}  & Prof. Tullio Vardanega, Prof. Riccardo Cardin, Eggon, BugBusters\\
\textbf{Versione} & 0.0.1\\
\end{tabularx}
\end{tcolorbox}
\end{center}

\vspace{0.5cm}

\begin{center}
\begin{tcolorbox}[colback=secondaryblue!10,colframe=secondaryblue,width=0.9\textwidth,arc=3mm,boxrule=0.8pt,title={\bfseries Abstract}]
Il presente verbale documenta la riunione esterna del 30 marzo 2026 con i referenti di Eggon, durante la quale il team ha illustrato lo stato di avanzamento dello sviluppo del progetto. 
L'azienda ha espresso parere positivo e ha formalmente approvato il prodotto realizzato. Sono stati mostrati il modulo di generazione post (AIAssistant) e
il modulo di estrazione e classificazione documenti (AICopilot), con le relative funzionalità\G di interfaccia\G.
Si è discusso delle scelte tecnologiche adottate, in particolare dei modelli di intelligenza artificiale utilizzati,
e sono stati concordati i prossimi passi riguardanti la documentazione tecnica.
\end{tcolorbox}
\end{center}

\newpage
\tableofcontents
\newpage

\section{Informazioni generali}

\begin{itemize}
    \item \textbf{Tipo riunione:} Esterna
    \item \textbf{Piattaforma:} Google Meet\G
    \item \textbf{Data:} 30/03/2026
    \item \textbf{Orario di inizio:} 14:00
    \item \textbf{Orario di fine:} 14:20
    \item \textbf{Presenti:}
    \begin{itemize}[leftmargin=1.5em, itemsep=3pt, label={\rule[0.5ex]{0.4em}{0.4em}}]
        \item Alberto Autiero
        \item Marco Favero
        \item Alberto Pignat
        \item Marco Piro
        \item Linor Sadè
        \item Leonardo Salviato
        \item Luca Slongo
    \end{itemize}
    \item \textbf{Assenti:}
    \begin{itemize}[leftmargin=1.5em, itemsep=3pt, label={\rule[0.5ex]{0.4em}{0.4em}}]
        \item Nessuno
    \end{itemize}
    \item \textbf{Presenti Esterni:}
    \begin{itemize}[leftmargin=1.5em, itemsep=3pt, label={\rule[0.5ex]{0.4em}{0.4em}}]
      \item Gianpaolo Ferrarin
      \item Luca Iuzzolino
    \end{itemize}
\end{itemize}

\section{Ordine del giorno}

\begin{enumerate}
    \item \label{itm:gen} Dimostrazione del modulo di generazione post (AI Assistant).
    \item \label{itm:ext} Dimostrazione del modulo di estrazione e classificazione documenti (AI Copilot).
    \item \label{itm:tech} Scelte tecnologiche: modelli di intelligenza artificiale adottati.
    \item \label{itm:next} Prossimi passi e documentazione tecnica.
\end{enumerate}

\newpage
\section{Svolgimento}

\subsection*{\ref{itm:gen}. Dimostrazione del modulo di generazione post (AI Assistant)}

Il team ha illustrato il funzionamento del generatore di post. Le aziende vengono gestite tramite database\G. Per ciascuna azienda è possibile definire uno stile e un tono comunicativo personalizzato. L'utente può inserire un prompt\G e avviare la generazione, ottenendo un post accompagnato da un'immagine. Sono state mostrate le seguenti funzionalità\G:

\begin{itemize}
    \item \textbf{Valutazione e modifica:} Dopo la generazione è possibile visualizzare i
    parametri utilizzati, modificarli, cambiare l'immagine e redigere il testo nell'editor
    integrato.
    \item \textbf{Salvataggio e storico:} Il post non viene inserito nello storico finché
    non viene esplicitamente salvato. La funzione \textit{duplica} permette di ripartire
    con gli stessi parametri, mentre \textit{rigenera} avvia una nuova generazione senza
    salvare il precedente risultato.
    \item \textbf{Dashboard:} La dashboard\G mostra sia i post salvati che quelli generati
    ma non salvati, con indicatori quali il rating medio e il numero di rigenerazioni
    immediate. È possibile filtrare i dati per periodo.
\end{itemize}

\subsection*{\ref{itm:ext}. Dimostrazione del modulo di estrazione e classificazione documenti (AI Copilot)}

Il team ha illustrato il flusso di elaborazione dei documenti caricati nel sistema. Il
modulo è compatibile con PDF\G, CSV e file immagine (PNG, JPEG). Il processo si articola
nelle seguenti fasi:

\begin{itemize}
    \item \textbf{Splitting e analisi:} Il sistema suddivide i documenti, assegna un ID
    univoco a ciascun segmento (composto da padre e figlio) e avvia l'analisi tramite
    modello di intelligenza artificiale. I campi estratti includono: destinatario, reparto,
    azienda e categoria.
    \item \textbf{Confidenza e intervento umano:} Per ogni campo estratto viene calcolata
    una percentuale di confidenza. In assenza di un match nel database\G, il campo viene
    comunque mostrato con l'indicazione di nessun match. L'utente può correggere manualmente
    i campi estratti (Human in the Loop\G); ogni modifica fa ripartire l'analisi del documento.
    \item \textbf{Invio e programmazione:} Il documento può essere inviato immediatamente
    oppure programmato per una data successiva. Lo stato del
    documento si aggiorna di conseguenza.
    \item \textbf{Gestione duplicati:} Se viene caricato un documento già presente nel sistema,
    questo viene automaticamente riconosciuto e mostrato all'utente.
    \item \textbf{Eliminazione:} L'eliminazione di un documento rimuove sia il documento padre
    che tutti i figli associati.
    \item \textbf{Dashboard:} La dashboard\G riporta la percentuale di confidenza media dei
    documenti analizzati e il numero di interventi manuali effettuati sulla categoria.
\end{itemize}

\subsection*{\ref{itm:tech}. Scelte tecnologiche: modelli di intelligenza artificiale adottati}

Il team ha illustrato le scelte sui modelli di intelligenza artificiale utilizzati:

\begin{itemize}
    \item \textbf{Generazione immagini:} È stato mantenuto il modello Canvas (versione 1),
    ritenuto conveniente in termini di costo rispetto alle alternative multimodali di recente
    introduzione.
    \item \textbf{Generazione testo:} Dopo una fase di valutazione dei modelli AWS Nova (Lite
    V1 e Pro V1), il team ha adottato Nova Lite V2, rilasciato a fine 2025. Questo modello
    offre prestazioni superiori rispetto al Pro della generazione precedente con un costo
    inferiore di circa il 40--45\%. Il modello rispetta meglio i prompt\G e restituisce il
    testo già formattato. È impiegato anche nella fase di splitting e analisi OCR\G dei documenti.
\end{itemize}

I referenti di Eggon hanno espresso pieno apprezzamento per il prodotto realizzato, 
approvandolo formalmente e lodando l’ottimo lavoro svolto dal team.

\vspace{1.0cm}
\noindent\textbf{\Large Firma per approvazione MVP}\\[0.4cm]
\underline{\hspace{8cm}}

\subsection*{\ref{itm:next}. Prossimi passi e documentazione tecnica}

Il team sta ultimando i documenti di Specifica Tecnica\G e Manuale Utente\G. Una volta conclusi, è stato concordato di condividerli con i referenti di Eggon non appena disponibili, per poi organizzare un incontro dedicato alla loro revisione. Contestualmente, il team intende candidarsi formalmente per il raggiungimento della Product Baseline (PB) con i docenti.

\section{Tabella delle decisioni e azioni}
\renewcommand{\arraystretch}{1.5}

% Impostazione globale del colore delle righe (fuori da tabularx)
\arrayrulecolor{primaryblue}
\sloppy
\begin{tabularx}{\textwidth}{
    |>{\raggedright\arraybackslash}p{3cm}|
    >{\raggedright\arraybackslash}X|
    >{\raggedright\arraybackslash}p{3cm}|
}
\hline
\rowcolor{primaryblue!40}
\textbf{\color{white} ID Decisione} & \textbf{\color{white} Descrizione} & \textbf{\color{white} Incaricato} \\
\hline
\DecisionRow{\href{https://github.com/BugBustersUnipd/DocumentazioneSWE/issues/224}{PB25}}{Stesura del verbale\G}{Alberto Autiero}
\DecisionRow{\href{https://github.com/BugBustersUnipd/DocumentazioneSWE/issues/225}{PB26}}{Verifica\G del verbale\G}{Alberto Pignat}
\DecisionRow{\href{https://github.com/BugBustersUnipd/DocumentazioneSWE/issues/226}{PB27}}{Portare a conclusione la documentazione propedeutica alla candidatura della Baseline PB}{Marco Favero, Linor Sadè, Leonardo Salviato}
\DecisionRow{\href{https://github.com/BugBustersUnipd/DocumentazioneSWE/issues/227}{PB28}}{Concludere la redazione del Piano di Qualifica e i test per la parte front-end dell'MVP}{Marco Piro}
\DecisionRow{\href{https://github.com/BugBustersUnipd/DocumentazioneSWE/issues/228}{PB29}}{Concludere la parte dei test per il back-end}{Luca Slongo}
\end{tabularx}
\fussy

\section{Esito Riunione}

Tutti i punti all'ordine del giorno sono stati discussi con esito positivo. 

Il gruppo BugBusters, preso atto dell'approvazione, procederà con la stesura della documentazione tecnica mancante e si candiderà ufficialmente per il raggiungimento della Product Baseline (PB) con i docenti.

Si ringraziano Gianpaolo Ferrarin e Luca Iuzzolino per la disponibilità e il supporto.

\vspace{1.5cm}
\noindent\textbf{\Large Data}\\[0.4cm]
\underline{\hspace{4cm}}

\vspace{1.5cm}
\noindent\textbf{\Large Firma}\\[0.8cm]
\underline{\hspace{6cm}} \\[0.2cm]
\vfill
\begin{center}
    {\small\color{darkgray} Documento redatto e approvato dal gruppo BugBusters.}
\end{center}

\end{document}